\subsection{Z-function: prefijos repetidos}\label{rec:z}

\codigo{z[i]} es el largo del prefijo más largo de \codigo{s} que coincide
con \codigo{s[i..]}. Para buscar un patrón se usa
\texttt{t = patron + '\#' + texto}.

\tablafunc{
\codigo{zfun(s)} & arreglo \codigo{z} del string \codigo{s} & $O(n)$ \\
\codigo{zsearch(t, p)} & posiciones de \codigo{p} en \codigo{t} & $O(n+m)$ \\
}

\copiar{el bloque \texttt{COPIA Z-FUNCTION}; \codigo{zsearch} ya separa con \texttt{'\#'}.}
\lstinputlisting{codigo/z_function.cpp}

El separador \texttt{'\#'} debe ser un carácter que no aparezca en las
strings. Si puede aparecer cualquiera, usa KMP.